% Part: lambda-calculus
% Chapter: church-rosser
% Section: parallel-beta-eta-reduction

\documentclass[../../../include/open-logic-section]{subfiles}

\begin{document}

\olfileid{lam}{cr}{pbe}

\olsection{并行 $\beta\eta$-归约}

本节证明并行 $\beta\eta$-归约的 Church-Rosser 性质，这是与 $\beta\eta$-归约相对应的并行归约概念。

\begin{defn}[并行 $\beta\eta$-归约，$\beredpar$] \ollabel{defn:beredpar}
  \emph{并行 $\beta\eta$-归约}（$\beredpar$）在项上归纳定义如下：
  \begin{enumerate}
    \item \ollabel{defn:beredpar1} $x \beredpar x$。
    \item \ollabel{defn:beredpar2} 若 $N \xrightarrow{\beta} N'$，则 $\lambd[x][N] \beredpar
      \lambd[x][N']$。
    \item \ollabel{defn:beredpar3} 若 $P \beredpar P'$ 且 $Q \beredpar Q'$，则 $PQ \beredpar
      P'Q'$。
    \item \ollabel{defn:beredpar4} 若 $N \beredpar N'$ 且 $Q \beredpar Q'$，则
      $(\lambd[x][N])Q \beredpar \Subst{N'}{Q'}{x}$。
    \item \ollabel{defn:beredpar5} 若 $N \beredpar N'$，则 $\lambd[x][Nx]
      \beredpar N'$，前提是 $x \notin FV(N)$。
  \end{enumerate}
\end{defn}

\begin{thm}\ollabel{thm:refl}
  $M \beredpar M$。
\end{thm}

\begin{proof}
  练习。
\end{proof}

\begin{prob}
  证明 \olref[lam][cr][pbe]{thm:refl}。
\end{prob}

\begin{defn}[$\beta\eta$-完全发展]\ollabel{defn:becd}
  $M$ 的\emph{$\beta\eta$-完全发展}~$\becd{M}$ 定义如下：
  \begin{align}
    \becd{x} &= x \ollabel{defn:becd1} \\
    \becd{(\lambd[x][N])} &= \lambd[x][\becd{N}] \ollabel{defn:becd2}\\
    \becd{(PQ)} &= \becd{P}\becd{Q} && \text{若 $P$ 不是 $\lambd$-抽象} 
    \ollabel{defn:becd3} \\
    \becd{((\lambd[x][N])Q)} &= \Subst{\becd{N}}{\becd{Q}}{x}
                             \ollabel{defn:becd4} \\
    \becd{(\lambd[x][Nx])} &= \becd{N} \ollabel{defn:becd5} & \text{若 $x
                                                      \notin FV(N)$}
  \end{align}
\end{defn}

\begin{lem}\ollabel{lem:comp}
  若 $M \beredpar M'$ 且 $R \beredpar R'$，则 $\Subst{M}{R}{y}
  \beredpar \Subst{M'}{R'}{y}$。
\end{lem}

\begin{proof}
  对 $M \beredpar M'$ 的推导进行归纳。

  前四种情况与 \olref[pb]{lem:comp} 中的完全相同。
  若最后一条规则是 \olref{defn:beredpar5}，则存在 $x$ 和 $N'$，使得 $M$ 是
  $\lambd[x][Nx]$，$M'$ 是 $N'$，其中 $x \notin
  FV(N)$ 且 $N \beredpar N'$。我们要证
  $\Subst{(\lambd[x][Nx])}{R}{y} \beredpar \Subst{N'}{R'}{y}$，即
  $\lambd[x][\Subst{N}{R}{y} x] \beredpar \Subst{N'}{R'}{y}$。这可由
  \olref{defn:beredpar}\olref{defn:beredpar5} 及归纳假设得出。
\end{proof}

\begin{lem}\ollabel{lem:cont}
  若 $M \beredpar M'$，则 $M' \beredpar \becd{M}$。
\end{lem}

\begin{proof}
  对 $M \beredpar M'$ 的推导进行归纳。

  前四种情况与 \olref[pb]{lem:cont} 中的类似。若
  最后一条规则是 \olref{defn:beredpar5}，则存在 $x$、$N$ 和 $N'$，使得 $M$ 是 $\lambd[x][Nx]$，
  $M'$ 是 $N'$，其中 $x \notin FV(N)$ 且 $N
  \beredpar N'$。我们要证 $N' \beredpar
  \becd{(\lambd[x][Nx])}$，即 $N' \beredpar \becd{N}$，这由归纳假设立得。
\end{proof}

\begin{thm}\ollabel{thm:cr}
  $\beredpar$ 具有 Church-Rosser 性质。
\end{thm}

\begin{proof}
  由 \olref{lem:cont} 直接可得。
\end{proof}

\end{document}